\bookchapter{Capitolo 26: Le frontiere tecnologiche dei telefoni attuali}

\chapterintro{Dopo aver parlato della “fine della vita” dei telefoni usati, riportiamo lo sguardo al presente. Il periodo 2024-2026 è una fase di esplosione concentrata delle tecnologie telefonistiche: l'IA è passata dal cloud alla tasca, la comunicazione satellitare è scesa dai top di gamma alla fascia media, gli schermi pieghevoli sono passati da “novità da provare” a “uso quotidiano”, e la tecnologia delle batterie ha conosciuto una svolta a livello di materiale senza precedenti da anni. In questo capitolo esaminiamo queste frontiere tecnologiche in corso, per aiutarti a decidere “se cambiare telefono ora” e “quale telefono scegliere”.}

\booksection{Telefoni IA: slogan di marketing o vera rivoluzione?}

Dal 2024, quasi tutti i produttori Android hanno lanciato sui loro top di gamma lo slogan “telefono IA”, e Apple ha introdotto Apple Intelligence in iOS 18. Ma cosa rende davvero un telefono un “telefono IA”? L'industria ha consolidato tre criteri di valutazione:

\begin{itemize}
  \item \textbf{Grandi modelli on-device}: il telefono è in grado di eseguire localmente modelli linguistici con circa 7B di parametri (come le versioni leggere di Qwen, Llama, DeepSeek), completando riassunti di testo, traduzioni, generazione di testi senza dipendere dal cloud. Metodo di verifica: se in modalità aereo le funzioni IA sono ancora disponibili, si tratta di modelli on-device.
  \item \textbf{Potenza di calcolo NPU (TOPS)}: la potenza della NPU (Neural Processing Unit) si misura in TOPS (trilioni di operazioni al secondo). I chip di punta del 2024 (come Snapdragon 8 Gen3, Dimensity 9300) hanno una potenza NPU tra 40 e 70 TOPS; le nuove generazioni 2025-2026 superano diffusamente i 100 TOPS, la soglia hardware per eseguire grandi modelli on-device.
  \item \textbf{Scenari AI Agent}: non solo "domande e risposte", ma anche la capacità di eseguire catene di task tra App diverse — ad esempio "organizza le foto della riunione di ieri in un album e condividilo con Mario", e il sistema richiama automaticamente Galleria, File, WeChat e altre App per completare l'operazione. Questo è il tratto distintivo che separa un telefono IA da un "telefono con assistente IA".
\end{itemize}

\booksection{Confronto delle strategie IA dei produttori: chi fa sul serio, chi si limita a etichettare?}

Le strategie IA dei vari brand differiscono notevolmente, e vale la pena prestare attenzione all'acquisto:

\begin{itemize}
  \item \textbf{Apple Apple Intelligence (iOS 18+)}: si concentra sull'integrazione di app a livello di sistema — riassunti email, priorità notifiche, revisione scrittura, fotoritocco IA in Galleria. Prediligendo modelli on-device, i task complessi vengono inviati al Private Cloud Compute. Vantaggi: protezione della privacy e fluidità di sistema; svantaggi: lento ad arrivare per gli scenari in cinese, alcune funzioni supportate solo progressivamente dal 2025.
  \item \textbf{Google Gemini Nano}: integrato nei Pixel 8/9, supporta "Cerca cerchiando", Magic Compose, riassunto registrazioni. Gemini Nano è un vero modello piccolo on-device, amichevole verso l'ecosistema open source, ed è l'implementazione di riferimento per i telefoni IA Android.
  \item \textbf{Modello Pangu di Huawei}: integrato nell'assistente Xiaoyi di HarmonyOS, supporta generazione di testi on-device, fotoritocco IA, rimozione oggetti IA (eliminare i passanti), riassunto chiamate. Con HarmonyOS NEXT 5.0 l'integrazione a livello di sistema è stata ulteriormente approfondita, fondendo le capacità IA con il sistema.
  \item \textbf{Modello Surge di Xiaomi (MiLM)}: integrato in HyperOS, supporta ritratti IA, riassunti IA, pittore IA, super-risoluzione. Xiaomi segue una strategia "on-device + cloud collaborativo": modelli leggeri on-device, task pesanti sul cloud.
  \item \textbf{OPPO (modello Andes) e vivo (modello Blue Heart)}: integrazione profonda in ColorOS 15 e OriginOS 5, con focus su rimozione oggetti IA, riassunto documenti IA, Q\&A su immagini IA, ricerca foto in linguaggio naturale ("trova le foto scattate al mare l'anno scorso").
\end{itemize}

\textbf{Consigli d'acquisto}: se usi spesso l'IA per elaborare documenti di lavoro o organizzare informazioni, privilegia dispositivi con forte potenza di calcolo on-device e integrazione sistemica delle funzioni IA; se usi solo occasionalmente il fotoritocco IA, i telefoni di fascia media sono già sufficienti.

\booksection{Comunicazione satellitare: da “novità” di punta a standard di fascia media}

La comunicazione satellitare è una delle tecnologie più sostanziose discese in fascia media degli ultimi due anni. Quando nel 2022 Huawei Mate 60 e iPhone 14 l'hanno lanciata, era ancora esclusiva dei top di gamma; nel 2025 è già scesa alla fascia 2000-3000 yuan.

\begin{itemize}
  \item \textbf{Messaggi brevi Beidou in fascia media}: la funzione di messaggi brevi del satellite Beidou (può inviare SMS fino a 1000 caratteri cinesi + posizione) è stata portata su modelli di fascia media come Redmi, realme, iQOO, senza canone mensile, gratuita. Adatta a escursionismo, viaggi in auto e simili.
  \item \textbf{Chiamate via satellite Tiantong}: Huawei Mate 60/70, Honor Magic6 e altri supportano la chiamata satellitare Tiantong-1, con possibilità di chiamare/ricevere telefonate satellitari e inviare messaggi satellitari. Richiede l'attivazione di un pacchetto satellitare dell'operatore (circa 10-200 yuan al mese), adatto a chi entra spesso in zone disabitate.
  \item \textbf{Messaggi satellitari iPhone (SOS)}: iPhone 14 e successivi supportano SOS via satellite, possono inviare messaggi di testo ai servizi di emergenza e condividere la posizione. Servizio gratuito per 2 anni, rinnovabile a pagamento in alcune regioni.
  \item \textbf{Consiglio pratico}: gli utenti urbani non ne hanno quasi mai bisogno, ma per gli appassionati di outdoor e chi viaggia spesso per lavoro in zone remote è fortemente raccomandato scegliere modelli con funzioni satellitari — nei momenti cruciali può salvare la vita.
\end{itemize}

\booksection{Maturità dell'ecosistema degli schermi pieghevoli: si può usare come telefono principale?}

Gli schermi pieghevoli sono passati da “novità tecnologica” a “pronti all'uso quotidiano”. Per valutare la maturità di un pieghevole, si guardano tre aspetti: durabilità hardware, adattamento delle App, ecosistema di interazione.

\begin{itemize}
  \item \textbf{Durabilità hardware}: la vita utile delle cerniere supera diffusamente le 500.000 aperture (100 al giorno per 13 anni), le cerniere a goccia riducono la piega dello schermo, la resistenza all'acqua IPX8 è ormai standard. La combinazione di vetro UTG + pellicola protettiva per lo schermo interno rende lo schermo meno "delicato".
  \item \textbf{Adattamento delle App}: le App principali (WeChat, Taobao, Douyin, Bilibili, Office, WPS) sono state adattate per i grandi schermi, ma circa il 30\% delle App di coda non è ancora adattato, con visualizzazione stirata o con bordi neri. Si consiglia di provare le App che usi più spesso prima dell'acquisto.
  \item \textbf{Interazione in stato di sospensione}: piegando a 90°-120°, alcune App supportano la modalità sospesa con "contenuto nella metà superiore, comandi nella metà inferiore", ideale per videocall, selfie, visione di video. Il selfie sospeso della fotocamera è un'esperienza esclusiva dei pieghevoli.
  \item \textbf{Split-screen e multi-tasking}: sul grande schermo si possono aprire due App in split-screen, o anche tre App in finestre piccole contemporaneamente. La logica di split-screen di HyperOS, ColorOS, OriginOS è ormai abbastanza matura, ma lo split-screen di iOS (in stile iPadOS) non è ancora disponibile su iPhone.
  \item \textbf{Prezzo}: i pieghevoli orizzontali (a libro) sono scesi dalla fascia oltre 10.000 yuan a 5.000-7.000 yuan; i verticali (a conchiglia) sono scesi a 4.000-5.000 yuan. Nel 2025 alcuni produttori hanno lanciato modelli verticali a 3.000 yuan, portando davvero i pieghevoli al mercato di massa.
\end{itemize}

\textbf{Consigli d'acquisto}: se sei un utente intensivo di lettura, documenti e video, lo schermo pieghevole orizzontale vale la pena; se cerchi portabilità e status sociale, quello verticale è più adatto. Si consiglia di provarlo in negozio per almeno 30 minuti prima di decidere.

\booksection{6G e Wi-Fi 7: velocità e confini della connettività di nuova generazione}

Anche se il 6G non è ancora davvero arrivato, le tecnologie correlate sono già in cammino:

\begin{itemize}
  \item \textbf{Progressi del 6G}: il 3GPP prevede di completare lo standard 6G entro il 2027, con commercializzazione intorno al 2030. La velocità teorica di picco del 6G può raggiungere 1 Tbps (50 volte il 5G), con latenza fino a 0,1 ms, e supporto all'"integrazione comunicazione-sensing" (communication + sensing): il telefono potrebbe persino fungere da radar a bassa precisione. Attualmente i produttori sono ancora in fase di verifica dei prototipi, lontano dai prodotti consumer.
  \item \textbf{Wi-Fi 7 (802.11be)}: dal 2024 è standard sui top di gamma, velocità teorica di picco fino a 46 Gbps (4 volte il Wi-Fi 6), supporta banda 320MHz, 4K-QAM, MLO (Multi-Link Operation). Nelle prove con router Wi-Fi 7, i top di gamme raggiungono 2-5 Gbps, 2-3 volte più veloci del Wi-Fi 6E. Se passi a un telefono Wi-Fi 7, ricorda di abbinare un router Wi-Fi 7 per sfruttarne tutte le prestazioni.
  \item \textbf{Bluetooth 5.4 e LE Audio}: i nuovi telefoni 2024-2026 supportano diffusamente Bluetooth 5.4 e LE Audio (audio a basso consumo), possono collegare simultaneamente due paia di cuffie per ascoltare la stessa canzone, e Auracast permette a più persone in luoghi pubblici di ascoltare in sincrono la stessa sorgente audio.
\end{itemize}

\booksection{Batterie al silicio-carbonio: l'autonomia conosce finalmente una svolta a livello di materiale}

L'autonomia dei telefoni è stata criticata per anni come “una carica al giorno”; dal 2024 la diffusione della tecnologia delle batterie ad anodo silicio-carbonio ha portato un miglioramento sostanziale.

\begin{itemize}
  \item \textbf{Principio tecnico}: la capacità teorica dell'anodo in grafite tradizionale è di circa 372 mAh/g, quella del silicio arriva a 4200 mAh/g (circa 10 volte), ma durante carica/scarica il silicio si espande del 300\%, tendendo a polverizzarsi. L'anodo silicio-carbonio, attraverso una struttura composita nano-silicio-carbonio, contiene l'espansione del silicio entro limiti accettabili, aumentando la densità di energia della batteria del 15\%-30\%.
  \item \textbf{Prestazioni reali}: i top di gamma 2024-2025 con anodo silicio-carbonio (come Xiaomi 14 Ultra, Honor Magic 6 Ultimate, vivo X100 Pro) integrano in corpi di 7-8 mm batterie da 5000-6000 mAh, circa 800-1000 mAh in più della generazione precedente; con uso intensivo durano un giorno e mezzo.
  \item \textbf{Ricarica rapida da 100W in abbinamento}: l'anodo silicio-carbonio è compatibile con la ricarica rapida; 100W/120W che caricano completamente in 17 minuti non sono più notizia. Nel 2025 alcuni modelli supportano 150W-240W, con carica all'80\% in 10 minuti.
  \item \textbf{Consigli d'acquisto}: se sei un utente intensivo, privilegia modelli post-2024 con la combinazione "anodo silicio-carbonio + ricarica rapida 100W", l'ansia da autonomia si ridurrà drasticamente.
\end{itemize}

\booksection{Panoramica del capitolo: una frase per ricordare queste frontiere}

\begin{itemize}
  \item Telefoni IA, tre criteri: grandi modelli on-device, potenza NPU in TOPS, scenari AI Agent tra App.
  \item Comunicazione satellitare: messaggi brevi Beidou in fascia media, chiamate Tiantong ancora esclusive dei top di gamma.
  \item Schermi pieghevoli: hardware maturo, adattamento App al 70\%, sospensione/split-screen come esperienze esclusive.
  \item Il 6G è ancora in cammino, il Wi-Fi 7 è arrivato, il Bluetooth LE Audio inizia a diffondersi.
  \item Anodo silicio-carbonio + ricarica 100W: doppia svolta per autonomia e ricarica.
\end{itemize}

\begin{tcbpractice}{Prova: prepara una "lista dei requisiti tecnologici" per il tuo prossimo telefono}
Prendi un foglio e, secondo tre fasce “indispensabile / bonus / non serve”, elenca le tue esigenze riguardo IA, satellite, schermo pieghevole, Wi-Fi 7, ricarica rapida, capacità batteria. Così, la prossima volta in negozio, potrai valutare rapidamente se i “punti di forza” consigliati dal commesso sono vere necessità o prodotti spinti. Ad esempio:

\begin{itemize}
  \item Indispensabile: batteria oltre 5000mAh, ricarica oltre 80W, NFC, resistenza IP68.
  \item Bonus: riassunti IA on-device, Wi-Fi 7, messaggi brevi Beidou.
  \item Non serve: schermo pieghevole (sono abituato al telefono piatto), chiamate satellitari (non vado in zone remote), ricarica estrema 240W.
\end{itemize}

\end{tcbpractice}

\begin{tcbsummary}{Riepilogo del capitolo}
\begin{itemize}
  \item Tre criteri per i telefoni IA: grandi modelli on-device, NPU TOPS, scenari AI Agent
  \item Strategie IA dei produttori: Apple focalizzata su privacy e integrazione, Google su on-device open source, Huawei/Xiaomi/OPPO/vivo ciascuno con i propri punti di forza
  \item Comunicazione satellitare: messaggi brevi Beidou in fascia media, chiamate Tiantong ancora esclusive dei top di gamma
  \item Schermi pieghevoli: hardware ormai maturo, adattamento App circa 70\%, stati di sospensione e split-screen come esperienze esclusive
  \item Standard 6G previsto per il 2027, Wi-Fi 7 ormai diffuso, Bluetooth LE Audio inizia a decollare
  \item Batterie al silicio-carbonio rendono 5000-6000mAh lo standard dei top di gamma, riducendo drasticamente l'ansia da autonomia
\end{itemize}

\end{tcbsummary}
